\answer
\begin{enumerate}

%%--- 章末問題 10.1 ---
\item[\ref{ex10.1}]
(a) 2段を直列に通るから，全体の効率は積になる。
\begin{equation*}
\eta = 0.95 \times 0.95 \approx 0.90
\end{equation*}

(b) 入力電力は$5/0.90 \approx 5.54$\,kWだから，損失はその差
\begin{equation*}
P_{\mathrm{loss}} = 5.54 - 5 \approx 0.54\,\mathrm{kW}
\end{equation*}
である。段数を重ねるほど効率は積で下がる ― 直接式が1段で変換したくなる理由の1つである。

%%--- 章末問題 10.2 ---
\item[\ref{ex10.2}]
$\alpha = 0$のときの実効値は$100$\,Vである。式\ref{eq10.3}の根号の中（＝電力比）を各角度で計算する。

$\alpha = 60^\circ = \pi/3$：$1 - 1/3 + \sin 120^\circ/(2\pi) = 0.667 + 0.138 = 0.805$。
実効値は$100\sqrt{0.805} \approx 89.7$\,V，電力は$0.80$倍。

$\alpha = 90^\circ = \pi/2$：$1 - 1/2 + 0 = 0.5$。
実効値は$70.7$\,V，電力は$0.50$倍（例題\ref{rei10.1}）。

$\alpha = 120^\circ = 2\pi/3$：$1 - 2/3 + \sin 240^\circ/(2\pi) = 0.333 - 0.138 = 0.196$。
実効値は$100\sqrt{0.196} \approx 44.2$\,V，電力は$0.20$倍。

%%--- 章末問題 10.3 ---
\item[\ref{ex10.3}]
導通するのは全周期$2\pi$のうち$\alpha \le \theta < \pi$だけだから
\begin{equation*}
V_{\mathrm{rms}}^2 = \frac{1}{2\pi} \int_\alpha^{\pi} V_0^2 \sin^2\theta\, d\theta
= \frac{V_0^2}{2\pi} \cdot \frac{1}{2}\left(\pi - \alpha + \frac{\sin 2\alpha}{2}\right)
= \frac{V_0^2}{4}\left(1 - \frac{\alpha}{\pi} + \frac{\sin 2\alpha}{2\pi}\right)
\end{equation*}
となり（積分は式\ref{eq10.2}と同じ），平方根を取れば与式を得る。
トライアック版と比べると，同じ点弧角で導通する期間が半分（正の半周期のみ）だから，
負荷へ渡るエネルギーも半分になる。
実効値の2乗は電力に比例するので，実効値は$1/\sqrt{2}$倍になるのである。

%%--- 章末問題 10.4 ---
\item[\ref{ex10.4}]
区間の長さは$\pi/3$だから
\begin{eqnarray}
\bar{V}_{\mathrm{out}}
&=& \frac{3}{\pi} \int_{\pi/3+\alpha}^{2\pi/3+\alpha} V_0 \sin\theta\, d\theta
= \frac{3 V_0}{\pi} \Bigl[ -\cos\theta \Bigr]_{\pi/3+\alpha}^{2\pi/3+\alpha} \nonumber\\
&=& \frac{3 V_0}{\pi} \left\{ \cos\left(\frac{\pi}{3}+\alpha\right)
- \cos\left(\frac{2\pi}{3}+\alpha\right) \right\} \nonumber
\end{eqnarray}
となる。和積の公式$\cos A - \cos B = -2\sin\frac{A+B}{2}\sin\frac{A-B}{2}$で
$A = \pi/3+\alpha$，$B = 2\pi/3+\alpha$とおくと
\begin{equation*}
\cos A - \cos B = -2\sin\left(\frac{\pi}{2}+\alpha\right)\sin\left(-\frac{\pi}{6}\right)
= 2\cos\alpha \cdot \frac{1}{2} = \cos\alpha
\end{equation*}
であり，式\ref{eq10.4}が得られる。平均出力は
$\alpha = 0^\circ$で$(3/\pi)V_0 \approx 0.95 V_0$，
$60^\circ$で$0.48 V_0$，$90^\circ$で$0$，
$120^\circ$で$-0.48 V_0$となり，負の電圧まで連続的に作れることがわかる。

%%--- 章末問題 10.5 ---
\item[\ref{ex10.5}]
(a) 線間電圧の振幅は$V_0 = \sqrt{2}\times400 \approx 566$\,V。式\ref{eq10.4}で$\alpha = 0$として
\begin{equation*}
\bar{V}_{\mathrm{max}} = \frac{3}{\pi}V_0 \approx 0.955\times566 \approx 540\,\mathrm{V}
\end{equation*}

(b) 目標は$\bar{V}_{\mathrm{out}}(t) = 300\sin(2\pi\times10\,t)$である。
式\ref{eq10.4}より$\cos\alpha(t) = 300\sin(2\pi\times10\,t)/540 = 0.556\sin(2\pi\times10\,t)$だから，
$\alpha$は$\cos\alpha = 0.556$の$56^\circ$（出力の山）から$\cos\alpha = -0.556$の$124^\circ$（出力の谷）までを$10$\,Hzで往復する。
出力の1周期$100$\,msの間に，入力の乗り継ぎ区間（$60^\circ$ごと，$1/300$\,s）は$30$回あり，
\ref{fig10.4}のようにこの30個の切れ端で正弦波1周期をなぞることになる。

(c) 出力線間電圧の実効値は入力の$\sqrt{3}/2$倍までだから，$400\times0.87 \approx 346$\,Vが上限である。
蓄えを持たない変換器は，どの瞬間も入力電圧の範囲の中でしか出力を作れない ― この制約が$0.87$倍の正体である（\ref{sec10.4}節）。

%%--- 章末問題 10.6 ---
\item[\ref{ex10.6}]
(a) 間接式は2段の積で$\eta = 0.96 \times 0.96 \approx 0.92$，
マトリックスコンバータは1段で$\eta = 0.96$である。

(b) 出力$10$\,kWに対する入力は，間接式が$10/0.92 \approx 10.85$\,kW（損失$0.85$\,kW），
マトリックスコンバータが$10/0.96 \approx 10.42$\,kW（損失$0.42$\,kW）。
1段で済む直接式は損失をおよそ半分にできる。

(c) 例：(1)~前段と後段をDCリンクで切り離して独立に設計・制御でき，
制御が単純で電源の乱れにも強い。
(2)~8章・9章の回路・部品・制御の実績と量産効果をそのまま使える
（ほかに，直接式は双方向スイッチの転流など制御が難しいこと，
出力電圧が入力の$0.87$倍までに制限されることなども挙げられる）。

%%--- 章末問題 10.7 ---
\item[\ref{ex10.7}]
(a) 式\ref{eq10.3}の根号の中（＝電力比）を$k$とおくと，
実効値は\rev{$V_{\mathrm{rms}} = 100\sqrt{k}$\,Vである（電源の実効値$141.4/\sqrt{2} = 100$\,V）}。各角度で$k$を計算して比べればよい。

$\alpha = 30^\circ$：$k = 1 - 1/6 + \sin 60^\circ/(2\pi) = 0.971$，実効値\rev{$98.5$}\,V（電力$0.97$倍）。

$\alpha = 60^\circ$：$0.805$，\rev{$89.7$}\,V（$0.80$倍）。

$\alpha = 90^\circ$：$0.500$，\rev{$70.7$}\,V（$0.50$倍）。

$\alpha = 120^\circ$：$0.196$，\rev{$44.3$}\,V（$0.20$倍）。

$\alpha = 150^\circ$：$0.029$，\rev{$17.0$}\,V（$0.03$倍）。

シミュレーションで読み取る実効値もこれらと一致する。
点弧角を遅らせるほど導通する区間が短くなり，実効値も電力も急速に下がっていく。

(b) $\alpha = 90^\circ$では根号の中がちょうど$1/2$になり，
実効値は電源（\rev{$100$}\,V）の$1/\sqrt{2}$倍＝\rev{$70.7$}\,Vとなる（例題\ref{rei10.1}）。
波形から求めた実効値もほぼこの値を示す。
半周期のうち後ろ半分だけを負荷に渡すので，
エネルギー（＝実効値の2乗）がちょうど半分になる点弧角である。

(c) $RL$負荷では，電圧が零になっても，インダクタに蓄えた磁気エネルギーが尽きるまで電流が流れ続ける。
このため消弧が電圧の零点より遅れ，抵抗負荷のような「電圧零＝電流零での自然消弧」が起こらない。
同じ点弧角でも導通区間が広がるので，実効値は抵抗負荷のときより大きくなる。
負荷が誘導性であるほど，この遅れは大きくなる。
\end{enumerate}
